% !TITLE 定风波·莫听穿林打叶声
% !AUTHOR 苏轼
% !DYNASTY 两宋
% !GENRE 词
% !ORDER 8

\begin{poem}
莫听穿林打叶声，何妨吟啸\textsuperscript{1}且徐行。
竹杖芒鞋\textsuperscript{2}轻胜马，谁怕？一蓑烟雨任平生\textsuperscript{3}。

料峭\textsuperscript{4}春风吹酒醒，微冷，山头斜照却相迎。
回首向来萧瑟\textsuperscript{5}处，归去，也无风雨也无晴\textsuperscript{6}。
\end{poem}

\begin{annotations}
\notetext{1}{吟啸：吟诗长啸。啸，撮口发出的长声，古人用它来表达情绪。雨中边走边唱，从容自在。}
\notetext{2}{芒鞋：草鞋。芒，一种草本植物，可编鞋。拄着竹杖、穿着草鞋，比骑马还轻快。}
\notetext{3}{一蓑烟雨任平生：披一领蓑衣，在烟雨中度过一生也无所谓。蓑（suō），用草或棕编的雨衣。“任平生”三个字是全词的灵魂——不管什么风雨，都随它去吧。}
\notetext{4}{料峭：形容春寒。料峭（liào qiào），微风带着寒意。}
\notetext{5}{萧瑟：风吹草木的声音，指风吹雨打的凄凉处。}
\notetext{6}{也无风雨也无晴：既没有风雨，也没有晴天——风雨和晴天都是外境，内心超越了这些分别。苏轼用这句话表达了一种豁达：在内心平静的人那里，外在的风雨或晴天都不再重要。}
\end{annotations}

\begin{appreciation}
这首词是苏轼被贬到黄州时写的。他因为"乌台诗案"差点丢了性命，被贬到黄州，在城东的坡地上种田，给自己取了个号叫"东坡居士"。这首《定风波》记录的只是一次淋雨，却是他一辈子心态的缩影，也是豪放词里最洒脱的作品之一。

"莫听穿林打叶声。"雨点穿过树林打在叶子上，哗哗作响，他选择不听。"何妨吟啸且徐行"，一边吟诗长啸，一边慢慢走。别人都在躲雨，他反而走得不慌不忙。"竹杖芒鞋轻胜马"，竹杖、草鞋，是穷人出门的装备，他说比骑马还轻快。"谁怕？"这一句反问，是全词的胆气。"一蓑烟雨任平生"，披一领蓑衣，一生的风雨，都随它去吧。

下阕写雨停了。料峭春风吹醒酒意，身上微冷，可"山头斜照却相迎"——太阳从山头探出来了。回头看看刚才风雨交加的地方，"也无风雨也无晴"。风雨和晴天都是外头的事，心里不挂碍，就都伤不着他。

书里《生年不满百》劝人"为乐当及时"，苏轼比它又进了一步：乐不乐都在当下，雨来了就淋着，雨走了就晒着，不抱怨，也不狂喜。人在坏运气里最容易低头，苏轼却是在最坏的年头里写出了最好的词。这点不服输的劲儿，比他的词本身更难得。

哥哥常想，你要是学会了这首词，往后我就不怎么担心你了。你以后也会碰上自己的风雨——考试、毕业、离家的头一个晚上。哥哥没法替你挡，你记住这七个字就够用了：也无风雨也无晴。
\end{appreciation}
